\documentclass[10pt,aspectratio=169]{beamer}

\usepackage[utf8]{inputenc}
\usepackage[T1]{fontenc}
\usepackage{lmodern}
\usepackage{amsmath,amssymb}

\usetheme{Madrid}
\usecolortheme{whale}

\title[Szczegolna Teoria Wzglednosci]{Podstawy Szczegolnej Teorii Wzglednosci}
\subtitle{Czas, przestrzen i rownowaznosc masy oraz energii}
\author[Jan Kowalski]{Jan Kowalski}
\institute[Uniwersytet]{Instytut Fizyki Teoretycznej}
\date{\today}

\begin{document}

% Slajd 1: Tytulowy
\begin{frame}
    \titlepage
\end{frame}

% Slajd 2: Spis tresci
\begin{frame}{\frametitle{Plan Prezentacji}}
    \tableofcontents
\end{frame}

% Slajd 3: Postulaty
\section{Postulaty Einsteina}
\begin{frame}{\frametitle{Dwa Postulaty Szczegolnej Teorii Wzglednosci}}
    \begin{block}{I Postulat: Zasada Wzglednosci}
        Prawa fizyki sa jednakowe we wszystkich inercjalnych ukladach odniesienia. Zadne doswiadczenie fizyczne przeprowadzone wewnatrz ukladu nie pozwala stwierdzic, czy uklad ten spoczywa, czy porusza sie ruchem jednostajnym prostoliniowym.
    \end{block}

    \vspace{0.3cm}
    
    \begin{alertblock}{II Postulat: Niezmiennosc Predkosci Swiatla}
        Predkosc swiatla w prozni $c$ ma taka sama wartosc we wszystkich inercjalnych ukladach odniesienia:
        \[
            c \approx 299\,792\,458\ \text{m/s}
        \]
    \end{alertblock}
\end{frame}

% Slajd 4: Transformacja Lorentza (Dwie kolumny)
\section{Kinematyka Relatywistyczna}
\begin{frame}{\frametitle{Transformacja Lorentza}}
    \begin{columns}[T]
        \begin{column}{0.48\textwidth}
            \begin{block}{Rownania Transformacji}
                Przejscie z ukladu $S$ do $S'$:
                \begin{itemize}
                    \item $x' = \gamma (x - vt)$
                    \item $y' = y$
                    \item $z' = z$
                    \item $t' = \gamma \left(t - \frac{vx}{c^2}\right)$
                \end{itemize}
            \end{block}
        \end{column}
        
        \begin{column}{0.48\textwidth}
            \begin{exampleblock}{Czynnik Lorentza}
                Wspolczynnik skalujacy:
                \[
                    \gamma = \frac{1}{\sqrt{1 - \frac{v^2}{c^2}}}
                \]
                \begin{itemize}
                    \item Dla $v \ll c$: $\gamma \approx 1$
                    \item Dla $v \to c$: $\gamma \to \infty$
                \end{itemize}
            \end{exampleblock}
        \end{column}
    \end{columns}
\end{frame}

% Slajd 5: Dylatacja i Kontrakcja
\begin{frame}{\frametitle{Dylatacja Czasu i Skrocenie Lorentza}}
    \begin{block}{Dylatacja Czasu}
        Zegar w ruchu tyka wolniej z punktu widzenia obserwatora:
        \[
            \Delta t = \gamma \Delta t_0 = \frac{\Delta t_0}{\sqrt{1 - \frac{v^2}{c^2}}}
        \]
    \end{block}

    \begin{block}{Kontrakcja Dlugosci}
        Obiekty poruszajace sie ulegaja skroceniu w kierunku ruchu:
        \[
            L = \frac{L_0}{\gamma} = L_0 \sqrt{1 - \frac{v^2}{c^2}}
        \]
    \end{block}
\end{frame}

% Slajd 6: Czasoprzestrzen
\begin{frame}{\frametitle{Struktura Czasoprzestrzeni}}
    \begin{block}{Interwal Czasoprzestrzenny}
        Niezmiennik transformacji Lorentza miedzy dwoma zdarzeniami:
        \[
            (\Delta s)^2 = (c \Delta t)^2 - (\Delta x)^2 - (\Delta y)^2 - (\Delta z)^2 = \text{const}
        \]
    \end{block}

    \begin{alertblock}{Typy Interwalow}
        \begin{itemize}
            \item $(\Delta s)^2 > 0$: Czasopodobny (mozliwy zwiazek przyczynowy)
            \item $(\Delta s)^2 = 0$: Swietlny (polaczenie sygnalem swietlnym)
            \item $(\Delta s)^2 < 0$: Przestrzeniopodobny (brak zwiazku przyczynowego)
        \end{itemize}
    \end{alertblock}
\end{frame}

% Slajd 7: Dynamika Relatywistyczna
\section{Dynamika i Rownowaznosc Masy-Energii}
\begin{frame}{\frametitle{Ped i Energia Relatywistyczna}}
    \begin{enumerate}
        \item \textbf{Ped relatywistyczny:}
        \[
            \vec{p} = \gamma m \vec{v} = \frac{m \vec{v}}{\sqrt{1 - \frac{v^2}{c^2}}}
        \]

        \item \textbf{Calkowita energia relatywistyczna:}
        \[
            E = \gamma m c^2 = E_k + m c^2
        \]

        \item \textbf{Energia spoczynkowa ($v=0$):}
        \[
            E_0 = m c^2
        \]
    \end{enumerate}
\end{frame}

% Slajd 8: Relacja Energia-Ped
\begin{frame}{\frametitle{Relacja Energia--Ped}}
    \begin{block}{Zwiazek Energii i Pedu}
        Rownanie prawdziwe rowniez dla czastek bezmasowych ($m=0$, np. fotonow):
        \[
            E^2 = (p c)^2 + (m c^2)^2
        \]
    \end{block}

    \begin{exampleblock}{Dla fotona ($m=0$):}
        \[
            E = p c \implies p = \frac{E}{c} = \frac{h \nu}{c}
        \]
    \end{exampleblock}
\end{frame}

% Slajd 9: Zastosowania
\section{Weryfikacja i Zastosowania}
\begin{frame}{\frametitle{Dowody Doswiadczalne i Zastosowania}}
    \begin{exampleblock}{Potwierdzenie w Fizyce i Technice}
        \begin{itemize}
            \item \textbf{Czas zycia mionow:} Miony atmosferyczne docieraja do Ziemi dzieki dylatacji czasu.
            \item \textbf{Akceleratory czastek (LHC):} Przyśpieszanie czastek wymaga uwzglednienia wzrostu pedu relatywistycznego.
            \item \textbf{Systemy GPS:} Korekcja czasu w zegarach satelitarnych.
            \item \textbf{Reakcje jadrowe:} Zamiana masy w energie wedlug wzoru $\Delta E = \Delta m c^2$.
        \end{itemize}
    \end{exampleblock}
\end{frame}

% Slajd 10: Zakonczenie
\begin{frame}
    \centering
    \Huge \textbf{Dziekuje za uwage!}
    
    \vspace{1cm}
    
    \Large Czy sa jakies pytania?
\end{frame}

\end{document}